% VI. Резултати. СКЕЛЕТ. Тук няма нито едно число, докато измерването не е проведено.
% Празните клетки са \TODO съзнателно: попълва се от отчета на измерването, не от памет.
\section{Експериментални резултати и анализ}
\label{sec:results}

\subsection{Закъснение по двата канала}

Измерването сравнява закъснението от произвеждането на стойността до нейното изобразяване
при доставка по WebSocket~\cite{rfc6455} и при периодично запитване, при еднаква честота
на обновяване. Формата на резултата е дадена в таблицата по-долу
\tabref{tab:latency}. Отчитат се средна стойност и мярка за разсейване по условията от
Раздел~\ref{sec:method}.

\begin{table}[H]
\caption{Закъснение от произвеждане до изобразяване по вид на канала}
\label{tab:latency}
\centering
\fontsize{11}{13.2}\selectfont
\begin{tabular}{@{}L{4.5cm}ccc@{}}
\toprule
\textbf{Канал} & \textbf{Средно} & \textbf{Разсейване} & \textbf{Брой опити} \\
\midrule
WebSocket            & \TODO{} & \TODO{} & \TODO{} \\
Периодично запитване & \TODO{} & \TODO{} & \TODO{} \\
\bottomrule
\end{tabular}
\end{table}

\TODO{коментар на резултата: потвърждава ли се прогнозата от Раздел V, и ако не се
потвърждава, каква е установената причина}

\subsection{Сравнение на двата формата за обмен}

За един и същ пакет със стойности се измерват обемът на съобщението по мрежата и времето
за неговия разбор на клиента при сериализация в разширяемия език за разметка~\cite{xml10}
и в нотацията за обекти~\cite{rfc8259}. Формата на резултата е в таблицата по-долу
\tabref{tab:formats}.

\begin{table}[H]
\caption{Обем и време за разбор по формат на обмена}
\label{tab:formats}
\centering
\fontsize{11}{13.2}\selectfont
\begin{tabular}{@{}L{4.5cm}ccc@{}}
\toprule
\textbf{Формат} & \textbf{Обем} & \textbf{Време за разбор} & \textbf{Брой опити} \\
\midrule
XML  & \TODO{} & \TODO{} & \TODO{} \\
JSON & \TODO{} & \TODO{} & \TODO{} \\
\bottomrule
\end{tabular}
\end{table}

\TODO{коментар: доколко разликата в обема се запазва след компресия на транспортно ниво,
и има ли случай, в който по-обемният формат е за предпочитане}

\subsection{Поведение при нарастващ брой клиенти}

Измерва се броят доставени пакети за единица време и закъснението при нарастващ брой
едновременни връзки, докато системата престане да поддържа зададената честота на
обновяване. Границата е резултат от измерването, а не проектна цел.
\TODO{графика с брой връзки по хоризонталната ос и доставени пакети по вертикалната, с
легенда за разчитане на показателите}

\subsection{Обобщение на измерените резултати}

\TODO{обобщение: кои от трите твърдения от Раздел V се потвърждават, кои се опровергават,
и какви са установените ограничения на самото измерване}
